\chapter{Deep Hedging and Machine Learning in Pricing}\label{ch:ml:deep-hedging-and-machine-learning-in-pricing}

A one-month at-the-money call on a stock whose volatility follows Heston's model is worth 2.25 on a spot of 100, and one point
of implied volatility is worth 0.115 of it. Hedged every day to its Black--Scholes delta with a transaction cost of 5 basis
points, the hedge pays 0.123 in costs on average: more than a volatility point, for a trader whose edge is often less. A
network that learns when not to trade pays 0.099 and carries less risk. This chapter uses machine learning where
derivatives desks use it: to hedge under frictions that the textbook ignores (\omterm{def:ml:deep-hedging-and-machine-learning-in-pricing:dh}{deep hedging}), to replace a slow pricer by a
fast approximation (surrogates, and the differential learning that makes them accurate with few samples), and to calibrate
a model in one pass. Each is measured against the exact answer that Book~5's engines provide, and each has a place where it
breaks.

\section{Hedging under frictions as a learning problem}

Book~5 (chapter~4) measured the discrete hedging error of the Black--Scholes delta and (chapter~26) the hedging band that
transaction costs call for. With costs, the best hedge depends on the current holding (trading back to delta is not worth it
if delta has moved little), on the risk the trader is willing to carry, and on everything the model leaves out; there is no
closed form beyond asymptotic ones. The problem is a stochastic control problem (Book~4, chapter~9) whose objective is a risk
measure of the hedged P\&L.

\begin{definition}[Convex risk measure, entropic risk measure, indifference price]\label{def:ml:deep-hedging-and-machine-learning-in-pricing:risk}
A \emph{convex risk measure}\index{convex risk measure} $\rho$ assigns to a random P\&L $X$ the amount of cash that makes it
acceptable, and is monotone, cash-invariant ($\rho(X + c) = \rho(X) - c$) and convex (Föllmer and Schied, 2002); expected
shortfall (Book~6, chapter~21) is one. The \emph{entropic risk measure}\index{entropic risk measure} is $\rho(X) =
\lambda^{-1}\log\E[e^{-\lambda X}]$, the certainty equivalent of exponential utility with risk aversion $\lambda$. The
\emph{indifference price}\index{indifference price} of a claim is the premium that leaves the seller exactly as well off, by
the risk measure, as not selling, each with its best trading strategy: $p = \inf_\delta\rho(-Z + G_\delta) -
\inf_\delta\rho(G_\delta)$, where $G_\delta$ is the gain of trading strategy $\delta$ after costs. When the underlying's
price is a martingale, as in the chapter's simulation, not trading is best without the claim and the second term is zero.
\end{definition}

\section{Deep hedging}

\begin{definition}[Deep hedging]\label{def:ml:deep-hedging-and-machine-learning-in-pricing:dh}
\emph{Deep hedging}\index{deep hedging} represents the hedging strategy by a \omterm{def:ml:neural-networks-for-noisy-tabular-data:mlp}{neural network} that maps what is observable at
each date (time, prices, current holdings) to the new holdings, and trains it on simulated paths to minimise a \omterm{def:ml:deep-hedging-and-machine-learning-in-pricing:risk}{convex risk
measure} of the hedged P\&L after all frictions (Buehler, Gonon, Teichmann and Wood, 2019).
\end{definition}

The chapter's network (\cref{lst:ml:dh:band}) is the no-transaction band network of Imaki and co-authors (2021): from the
time to maturity and the log-moneyness it outputs a band, and the new holding is the old one clipped into it, so that not
trading is the default and the band's width can be read off. The P\&L of selling the call and hedging daily with the
underlying is computed on a batch of paths (\cref{lst:ml:dh:pnl}), the entropic risk with $\lambda = 1$ is the loss, and 1\,000
\omterm{def:ml:neural-networks-for-noisy-tabular-data:adam}{Adam} steps on batches of 2\,048 of 20\,000 training paths take about fifteen seconds on one core. Every \omterm{def:ml:reinforcement-learning-foundations:mdp}{policy} is then scored
on 20\,000 other paths.

\begin{omfigure}
\begin{tikzpicture}
\begin{axis}[width=11cm,height=6cm,xlabel={proportional cost (basis points)},ylabel={indifference price (entropic, $\lambda = 1$)},
  tick label style={font=\small},label style={font=\small},xmin=0,xmax=20,ymin=2.4,ymax=3.05,grid=major,grid style={gray!25},
  legend cell align=left,legend style={font=\footnotesize,at={(0.02,0.98)},anchor=north west,fill=white,draw=none},
  table/col sep=comma]
\addplot[omDef,very thick,mark=*] table[x=cost,y=deep]{figdata/ml/19-deep-hedging-and-machine-learning-in-pricing/indifference.csv};
\addplot[omProp,very thick,mark=square*] table[x=cost,y=ww]{figdata/ml/19-deep-hedging-and-machine-learning-in-pricing/indifference.csv};
\addplot[omThm,very thick,mark=triangle*] table[x=cost,y=bs]{figdata/ml/19-deep-hedging-and-machine-learning-in-pricing/indifference.csv};
\legend{deep hedge (band network),Whalley--Wilmott band,Black--Scholes delta}
\end{axis}
\end{tikzpicture}

\omcaption{The seller's \omterm{def:ml:deep-hedging-and-machine-learning-in-pricing:risk}{indifference price} of the one-month at-the-money call (its Heston value is 2.25) under each hedging
\omterm{def:ml:reinforcement-learning-foundations:mdp}{policy}, against the cost of trading the underlying. Data: \texttt{ml\_hedge.hedging}.}\label{fig:ml:dh:indiff}
\end{omfigure}

\Cref{fig:ml:dh:indiff} is the first half of the chapter's named result. Without costs the deep hedge's \omterm{def:ml:deep-hedging-and-machine-learning-in-pricing:risk}{indifference price} is
2.454 against 2.475 for the Black--Scholes delta: under Heston, spot and volatility move together ($\rho = -0.7$) and the
network learns a hedge ratio that uses it. With costs of 5, 10 and 20 basis points it rises to 2.558, 2.658 and 2.840; the
Black--Scholes delta hedge's to 2.608, 2.743 and 3.018, and the Whalley--Wilmott band's (Whalley and Wilmott, 1997) to 2.606,
2.712 and 2.910. \Cref{tab:ml:dh:ten} shows the other statistics at 10 basis points. Trained on expected shortfall instead, the
network lowers the expected shortfall (3.763 against 3.817) and raises the entropic risk (2.671 against 2.658): the risk
measure is part of the model, and each network is best by its own.

\begin{table}[htbp]
\centering\small
\begin{tabular}{@{}lcccc@{}}
\toprule
\omterm{def:ml:reinforcement-learning-foundations:mdp}{policy} & entropic risk & expected shortfall 95\% & mean cost & standard deviation\\
\midrule
deep hedge, entropic & 2.658 & 3.817 & 2.442 & 0.643\\
deep hedge, expected shortfall & 2.671 & 3.763 & 2.457 & 0.657\\
Whalley--Wilmott band & 2.712 & 4.044 & 2.424 & 0.703\\
Black--Scholes delta & 2.743 & 4.090 & 2.507 & 0.595\\
no hedge & 8.767 & 9.992 & 2.235 & 2.945\\
\bottomrule
\end{tabular}
\caption{Selling the call and hedging daily at 10 basis points of cost, on 20\,000 test paths: risk measures of the P\&L
before the premium, its mean loss and standard deviation. Data: \texttt{ml\_hedge.hedging}, \texttt{shortfall\_hedge}.}\label{tab:ml:dh:ten}
\end{table}

The second half of the named result is the band itself (\cref{fig:ml:dh:bands}). At the money half-way to maturity, the
learned band is 0.031 wide in delta at 5 basis points and 0.104 at 20; the Whalley--Wilmott band, $\pm(3cS\Gamma^2/2\lambda)^{1/3}$,
is 0.181 and 0.288. The asymptotic formula assumes a trader who watches continuously and trades at the band's edge; with one
decision a day, the position drifts between decisions and a wide band compounds the drift, so the network keeps a narrower
band at the money and above it, and a wider one below the money, where the option is nearly worthless to hedge. The formula gives the right idea and the wrong size
for this trader; the network finds the size for the trader's actual schedule and risk measure.

\begin{omfigure}
\begin{tikzpicture}
\begin{axis}[width=11cm,height=6cm,xlabel={spot (strike 100, half-way to maturity)},ylabel={holding (shares per option)},
  tick label style={font=\small},label style={font=\small},xmin=92,xmax=108,ymin=-0.05,ymax=1.0,grid=major,grid style={gray!25},
  legend cell align=left,legend style={font=\footnotesize,at={(0.02,0.98)},anchor=north west,fill=white,draw=none},
  table/col sep=comma]
\addplot[black,thick,dashed] table[x=spot,y=delta]{figdata/ml/19-deep-hedging-and-machine-learning-in-pricing/bands.csv};
\addplot[omProp,thick] table[x=spot,y=wwlo]{figdata/ml/19-deep-hedging-and-machine-learning-in-pricing/bands.csv};
\addplot[omDef,very thick] table[x=spot,y=deeplo]{figdata/ml/19-deep-hedging-and-machine-learning-in-pricing/bands.csv};
\addplot[omProp,thick,forget plot] table[x=spot,y=wwhi]{figdata/ml/19-deep-hedging-and-machine-learning-in-pricing/bands.csv};
\addplot[omDef,very thick,forget plot] table[x=spot,y=deephi]{figdata/ml/19-deep-hedging-and-machine-learning-in-pricing/bands.csv};
\legend{Black--Scholes delta,Whalley--Wilmott band,learned band}
\end{axis}
\end{tikzpicture}

\omcaption{No-trade bands at 20 basis points of cost, half-way to maturity: the holding is left alone inside a band and moved
to its nearer edge outside it. Data: \texttt{ml\_hedge} (\texttt{fig\_hedge.py}).}\label{fig:ml:dh:bands}
\end{omfigure}

\section{Surrogate pricers and differential learning}

\begin{definition}[Surrogate model, differential machine learning]\label{def:ml:deep-hedging-and-machine-learning-in-pricing:surrogate}
A \emph{surrogate model}\index{surrogate model} is a fast approximation of a slow function, here a pricer, fitted to its
outputs on sampled inputs. \emph{Differential machine learning}\index{differential machine learning} trains the surrogate on the
derivatives of the labels with respect to the inputs as well as on the labels, the network's own derivatives being computed by
automatic differentiation (Book~4, chapter~28); with Monte Carlo labels the derivative labels are pathwise (Huge and Savine,
2020).
\end{definition}

A risk system that prices a book under thousands of scenarios cannot afford a Fourier integral or a Monte Carlo run per
price. The chapter's surrogate learns the Heston price of a three-month call as a function of the spot (70 to 130) and the
initial variance (0.01 to 0.09) from the cheapest labels there are: one simulated path per sample, whose payoff is an unbiased
but very noisy price. The differential label is the pathwise derivative of the payoff with respect to the spot,
$\mathbf1\{S_T > K\}S_T/S_0$, free once the path exists. \Cref{tab:ml:dh:surr} scores both trainings against the Fourier price
on 400 test points (\cref{lst:ml:dh:surr} is the training loop).

\begin{table}[htbp]
\centering\small
\begin{tabular}{@{}rcccc@{}}
\toprule
& \multicolumn{2}{c}{price RMSE} & \multicolumn{2}{c}{delta RMSE}\\
\cmidrule(lr){2-3}\cmidrule(l){4-5}
training samples & standard & differential & standard & differential\\
\midrule
256 & 0.927 & 0.679 & 0.066 & 0.057\\
1\,024 & 0.436 & 0.218 & 0.059 & 0.021\\
4\,096 & 0.297 & 0.279 & 0.042 & 0.020\\
\bottomrule
\end{tabular}
\caption{A two-layer surrogate of the Heston call price learned from one-path Monte Carlo payoffs, with and without pathwise
delta labels: errors against the Fourier price and its finite-difference delta on 400 test points. Data:
\texttt{ml\_hedge.surrogates}.}\label{tab:ml:dh:surr}
\end{table}

With 1\,024 samples the derivative labels halve the price error and cut the delta error by nearly two-thirds; with 4\,096 the
prices are as good either way but only the differential network's deltas are. A surrogate's delta is what the hedger uses, and
a network fitted to prices alone has no reason to get its slope right.

\section{Calibration networks}

\begin{definition}[Deep calibration]\label{def:ml:deep-hedging-and-machine-learning-in-pricing:calib}
\emph{Deep calibration}\index{deep calibration} replaces the numerical optimisation of a model's parameters against market quotes
(Book~4, chapter~24) by a network trained on model-generated data: either the inverse map from quotes to parameters (Hernandez,
2016) or a fast forward map from parameters to quotes inside a standard optimiser (Horvath, Muguruza and Tomas, 2021).
\end{definition}

The chapter's network learns the inverse map from fifteen normalised call prices (five strikes, three maturities) to the five
Heston parameters, on 4\,000 parameter sets drawn uniformly in a box. On 200 \omterm{def:ml:why-financial-machine-learning-is-different:sets}{test sets}, its errors as a share of each
parameter's range are 1.7\% for the initial variance, 18.7\% for the mean-reversion speed, 9.7\% for the long-run variance,
9.0\% for the volatility of variance and 9.8\% for the correlation; the surfaces re-priced with its answers miss the true ones
by 13.2 basis points of spot on average. Speed and long-run variance are poorly identified by three maturities, and the network
is as uncertain about them as the data are. On five test surfaces its answer misses the implied volatilities by 0.46 points on
average, and Book~5's Levenberg--Marquardt calibrator, started from it or from a fixed guess, fits them exactly. The network is
a fast first guess and a sanity check, not a replacement, on clean data from the model it was trained on; on market data,
which no model fits exactly, it has never seen the residuals it will meet.

\section{Where learned pricing breaks}

Three failures recur. \emph{Extrapolation}: a surrogate or a calibration network outside its training box returns confident
numbers with no warning, exactly where a stressed market takes the book; train on boxes wider than any scenario and check
inputs against them. \emph{Arbitrage}: a surrogate's prices need not be monotone or convex in the strike, and its Greeks can
change sign where they should not; constrain the architecture or check the outputs. \emph{Model risk}: a deep hedge is optimal
for the simulator it was trained on, like chapter~17's agent; train it on several models or randomised parameters, and compare
it with the asymptotic band on paths it was not trained on.

\begin{method}[Machine learning on a derivatives desk]\label{met:ml:dh:method}
\begin{enumerate}
\item Keep the exact pricer as the reference; every surrogate is validated against it on a grid wider than the book's
scenarios, prices and Greeks both.
\item Train surrogates with differential labels when derivatives are cheap (pathwise or by automatic differentiation).
\item Use calibration networks as starting points and checks; the optimiser has the last word.
\item Train deep hedges on the desk's own frictions and risk measure, and compare with the closed-form band out of sample and
across models.
\end{enumerate}
\end{method}

\section{Tutorial: hedging with frictions}\label{tut:ml:deep-hedging-and-machine-learning-in-pricing:tut}

\begin{tutorial}
\textbf{Goal.} Deep-hedge the call at four cost levels and read off the \omterm{def:ml:deep-hedging-and-machine-learning-in-pricing:risk}{indifference prices} and bands; train the surrogate with
and without differential labels; fit the calibration network. \textbf{End state:}
\cref{fig:ml:dh:indiff,fig:ml:dh:bands}, \cref{tab:ml:dh:ten,tab:ml:dh:surr}.
\begin{tutsteps}
\item \textbf{The no-transaction band network.}
\omcode{code/firm/deephedge/firm_deephedge.py}{54}{69}{A network that outputs a band; the holding is clipped into it.}\label{lst:ml:dh:band}
\item \textbf{The hedged P\&L after costs.}
\omcode{code/firm/deephedge/firm_deephedge.py}{72}{86}{Selling a call and hedging it on a batch of paths.}\label{lst:ml:dh:pnl}
\item \textbf{Differential training.}
\omcode{code/firm/deephedge/firm_deephedge.py}{165}{186}{A surrogate trained on prices and pathwise derivatives.}\label{lst:ml:dh:surr}
\item \textbf{Run} \texttt{ml\_hedge.hedging()} (about a minute and a half on one core), \texttt{surrogates()},
\texttt{calibration()}, \texttt{lm\_check()} and \texttt{fig\_hedge.py}.
\end{tutsteps}
\textbf{What to change next.} Give the network the current variance as an input and add a variance swap as a second hedge;
train on randomised Heston parameters and test on another model.
\end{tutorial}

\section{Build: deep hedging and learned pricers}\label{bld:ml:deep-hedging-and-machine-learning-in-pricing:deephedge}

\begin{build}
\bfield{Purpose} Hedges and prices learned where they save time or cost, checked against Book~5's exact engines.
\bfield{Interface} \texttt{heston\_paths}; \texttt{BandNet}, \texttt{hedge\_pnl(policy, S, K, T, cost)}, \texttt{entropic},
\texttt{expected\_shortfall}, \texttt{train\_hedge(S, K, T, cost, risk, seed, steps, batch, lr)}; \texttt{bs\_policy},
\texttt{ww\_policy}, \texttt{no\_hedge}, \texttt{band}; \texttt{SurrogateNet}, \texttt{fit\_surrogate(X, y, dy,
differential, seed)}, \texttt{surrogate\_values}; \texttt{fit\_calibrator(vols, params, seed)}.
\bfield{Rules} Every learned price, Greek and hedge is scored against an exact reference on data it was not trained on;
training boxes and cost levels are recorded with the model.
\bfield{Acceptance tests} \texttt{code/firm/deephedge/tests/}: the P\&L of the delta hedge without costs matches the discrete
hedging error's scale; the entropic risk of a constant is minus the constant; the band network holds its position inside the
band; Whalley--Wilmott's band widens with cost; a differential surrogate of a known function recovers its slope; the
calibrator inverts a simple map.
\bfield{Stretch} A second hedging instrument; recurrent policies with the variance as a hidden state; arbitrage-free surrogate
architectures.
\end{build}

\begin{omsources}
\item H.~Buehler, L.~Gonon, J.~Teichmann and B.~Wood, ``Deep hedging'', \emph{Quantitative Finance} 19(8), 2019.
\item S.~Imaki, K.~Imajo, K.~Ito, K.~Minami and K.~Nakagawa, ``No-transaction band network'', arXiv:2103.01775, 2021.
\item A.~E.~Whalley and P.~Wilmott, ``An asymptotic analysis of an optimal hedging model for option pricing with transaction
costs'', \emph{Mathematical Finance} 7(3), 1997.
\item H.~Föllmer and A.~Schied, ``Convex measures of risk and trading constraints'', \emph{Finance and Stochastics} 6(4), 2002.
\item B.~Huge and A.~Savine, ``Differential machine learning'', arXiv:2005.02347, 2020.
\item B.~Horvath, A.~Muguruza and M.~Tomas, ``Deep learning volatility'', \emph{Quantitative Finance} 21(1), 2021.
\item A.~Hernandez, ``Model calibration with neural networks'', SSRN 2812140, 2016.
\end{omsources}

\section{Exercises}

\begin{exercise}[$\star$]\label{exo:ml:deep-hedging-and-machine-learning-in-pricing:1}
Show that the \omterm{def:ml:deep-hedging-and-machine-learning-in-pricing:risk}{entropic risk measure} is cash-invariant: $\rho(X + c) = \rho(X) - c$. What is $\rho$ of a constant?
\end{exercise}

\begin{exercise}[$\star$]\label{exo:ml:deep-hedging-and-machine-learning-in-pricing:2}
Compute the Whalley--Wilmott half-width at the money half-way to maturity for $c = 0.001$, $S = 100$, $\lambda = 1$ and a
Black--Scholes gamma of 0.0998.
\end{exercise}

\begin{exercise}[$\star$]\label{exo:ml:deep-hedging-and-machine-learning-in-pricing:3}
Why is the \omterm{def:ml:deep-hedging-and-machine-learning-in-pricing:risk}{indifference price} above the Heston value (2.25) even without costs?
\end{exercise}

\begin{exercise}[$\star\star$]\label{exo:ml:deep-hedging-and-machine-learning-in-pricing:4}
Why does the deep hedge beat the Black--Scholes delta when there are no costs at all?
\end{exercise}

\begin{exercise}[$\star\star$]\label{exo:ml:deep-hedging-and-machine-learning-in-pricing:5}
Why are the mean-reversion speed and the long-run variance hard for the calibration network, and what data would help?
\end{exercise}

\begin{exercise}[$\star\star$]\label{exo:ml:deep-hedging-and-machine-learning-in-pricing:6}
\emph{Find the flaw.} ``Our surrogate matches the pricer to 0.1\% on the \omterm{def:ml:why-financial-machine-learning-is-different:sets}{test set}, so we use it for the stress scenarios.''
\end{exercise}

\begin{exercise}[$\star\star\star$]\label{exo:ml:deep-hedging-and-machine-learning-in-pricing:7}
\emph{Coding.} Derive the pathwise delta label for a digital call ($\mathbf1\{S_T > K\}$). What goes wrong, and what does
Huge and Savine's approach do instead?
\end{exercise}

\begin{exercise}[$\star\star\star$]\label{exo:ml:deep-hedging-and-machine-learning-in-pricing:8}
Show that without costs and in a complete market (Black--Scholes, continuous trading) the entropic \omterm{def:ml:deep-hedging-and-machine-learning-in-pricing:risk}{indifference price} equals
the Black--Scholes price for any $\lambda$.
\end{exercise}

\section{Problem: Hedging with Frictions}

\begin{problem}[{Weekend problem --- when not to trade}]\label{pb:ml:deep-hedging-and-machine-learning-in-pricing:1}
The chapter's call, hedges, surrogate and calibrator.

\medskip\noindent\textbf{Part I --- The problem.}
\begin{enumerate}
\item What is sold, how is it hedged, and what does the network see?
\item What does daily delta hedging cost at 5 basis points, and how does that compare with the vega?
\item Why a band network rather than a network that outputs the holding directly?
\item What is the loss, and what does $\lambda = 1$ mean here?
\end{enumerate}

\medskip\noindent\textbf{Part II --- Results.}
\begin{enumerate}[resume]
\item Give the \omterm{def:ml:deep-hedging-and-machine-learning-in-pricing:risk}{indifference prices} of the three hedges at each cost level.
\item Compare the learned band with the Whalley--Wilmott band.
\item What does training on expected shortfall change?
\item Why does the Whalley--Wilmott hedge have the lowest mean cost at 10 basis points but not the lowest risk?
\end{enumerate}

\medskip\noindent\textbf{Part III --- Surrogates and calibration.}
\begin{enumerate}[resume]
\item What do differential labels buy, and when?
\item What errors does the calibration network make?
\item How should the network and the optimiser be combined?
\item Where would each learned pricer fail on a real desk?
\end{enumerate}

\medskip\noindent\textbf{Part IV --- The verdict.}
\begin{enumerate}[resume]
\item State the \emph{named result}: the \omterm{def:ml:deep-hedging-and-machine-learning-in-pricing:risk}{indifference price} as a function of the cost level, and the learned no-trade band's
width against the Whalley--Wilmott band.
\item Would you quote the option at the \omterm{def:ml:deep-hedging-and-machine-learning-in-pricing:risk}{indifference price}? Why or why not?
\item How would you validate a deep hedge before using it on a book?
\item What would change with a second hedging instrument?
\item How do you keep a surrogate honest when the model is recalibrated daily?
\item What is the risk of training the hedge on one Heston parameter set?
\item Where does automatic differentiation fit in this chapter?
\item In one sentence: what does a hedging network learn that a delta does not?
\end{enumerate}
\end{problem}

\section{Interview questions}

\begin{interviewq}[$\star$ \iqroles{researcher, trader}]\label{iq:ml:deep-hedging-and-machine-learning-in-pricing:1}
Why does delta hedging with transaction costs call for a no-trade band?
\end{interviewq}

\begin{interviewq}[$\star\star$ \iqroles{researcher}]\label{iq:ml:deep-hedging-and-machine-learning-in-pricing:2}
What is \omterm{def:ml:deep-hedging-and-machine-learning-in-pricing:dh}{deep hedging}, and what are its inputs, outputs and loss?
\end{interviewq}

\begin{interviewq}[$\star\star$ \iqroles{mle, researcher}]\label{iq:ml:deep-hedging-and-machine-learning-in-pricing:3}
Explain \omterm{def:ml:deep-hedging-and-machine-learning-in-pricing:surrogate}{differential machine learning} and why it helps with noisy Monte Carlo labels.
\end{interviewq}

\begin{interviewq}[$\star\star$ \iqroles{researcher}]\label{iq:ml:deep-hedging-and-machine-learning-in-pricing:4}
What is an \omterm{def:ml:deep-hedging-and-machine-learning-in-pricing:risk}{indifference price}, and how does it relate to a risk measure?
\end{interviewq}

\begin{interviewq}[$\star\star$ \iqroles{mle}]\label{iq:ml:deep-hedging-and-machine-learning-in-pricing:5}
A calibration network is fast. Why not use it alone?
\end{interviewq}

\begin{interviewq}[$\star\star\star$ \iqroles{researcher}]\label{iq:ml:deep-hedging-and-machine-learning-in-pricing:6}
Your surrogate's gamma turns negative for deep out-of-the-money strikes. What happened and how do you fix it?
\end{interviewq}
